\section{Discussion}
\label{sec:discussion}

Read-speech corpora for low-resource languages are built the way this one was:
a small pool of sentences recorded many times over by the few speakers who
volunteer -- here 2{,}067 sentences and 77 speakers. Any split that separates speakers while using all the recordings therefore
evaluates with complete text overlap, by construction rather than oversight:
the training portion must contain every sentence to contain every clip. We report this as a
property of the benchmark, and it is consequential: a development set whose
sentences the model has already heard from other speakers prefers systems that
remember text over systems that recognise speech.

The remedy is cheap: report the text-disjoint figure alongside the conventional
one. Where a corpus ships an official split disjoint in both respects, as
Common Voice does here at the cost of one recording per sentence, the lost
training data buys one honest measurement among the optimistic ones.

Decoding is the other place where the shape of the data decides the method.
Greedy CTC decoding commits to the most probable symbol at each 20\,ms frame
independently, with no notion of the word being formed, and the error profile
shows what that costs: on the baseline's out-of-domain set we measure 23.4\%
substitutions against almost no deletions or insertions -- the signature of
words that are nearly right. That invites a language model, and it must be over
characters: Quechua builds words by chaining suffixes, and 77\% of the distinct
word forms in the scripted sentence corpus occur exactly once, so a word-level
model would never have seen the correct form of most words it must judge.

% =====================================================================
